In this section, we study the behavior of the ramp loss model on some simple, artificial data.
These examples employ a ``full master problem'', meaning all relevant rules (columns) are considered.
Precisely, given a positive integer $D$, all possible rules with at most $D$ clauses satisfying the following conditions are considered:
\begin{itemize}
    \item A feature appears at most once in any rule.
    So a feature appears in at most one clause within a rule.
    \item A feature and their negated version do not appear together in any rule.
    This means a feature cannot appear in both a ``$\geq$''-clause and a ``$\leq$''-clause within a rule.
\end{itemize}
The examples are in two dimensions.
Recall that for $\beta = 0$, the ramp loss model is equivalent to the threshold model.

\begin{example}[Two clouds]\label{ex:two_clouds}
    \begin{figure}
        \includesvg[width=\textwidth]{gen/test_clouds/example_1.svg}
        \caption{Instance of \autoref{ex:two_clouds}. The example consists of a negative and a positive cloud of $100$ points each, following $\mathcal{N}((8, 10), 2)$ and $\mathcal{N}((12, 10), 2)$ respectively.}
        \label{fig:two_clouds}
    \end{figure}
    We construct two clouds of points according to multivariate normals.
    Precisely,
    \begin{itemize}
        \item A cloud of $100$ negative points drawn from $\text{Normal}((8, 10), 2)$,
        \item A cloud of $100$ positive points drawn from $\text{Normal}((12, 10), 2)$.
    \end{itemize}
    Hence, the cloud centers share $x_2$ but differ in $x_1$.
    An instance is depicted in \autoref{fig:two_clouds}
    Choosing $c=1$ as rule complexity, we solve the ramp loss model for the parameters
    \[C \in \{1,2\}, \quad \beta \in \{0, 0.5, 1, 2, 3, 4, 5\}, \quad D = \{2\}.\]

    \begin{figure}
        \includesvg[width=.8\textwidth]{gen/test_clouds/results/O_1_C_1_mw_0_additive.set.svg}
        \caption{Solution to an instance of \autoref{ex:two_clouds} computed with the ramp loss model, configured with $C=1$, $\beta=0$, $D=2$.}
        \label{fig:two_clouds_solution_simple}
    \end{figure}
    For $\beta=0,$ the solution contains only the rule
    $x_1 \geq 9.88$ regardless of $C$.
    So, for this example, an increase in $C$ does not cause the model to behave differently.
    This solution is depicted in \autoref{fig:two_clouds_solution_simple}.

    For $C=1$ fixed, the same solution is found for $\beta$ up until and including $\beta = 4$.
    For $C=1, \; \beta=5,$ the rule is $x_1 \geq 10.73$ instead.
    We conclude that for $C=1$, the ramp loss model finds a sensible solution for any of the tested $\beta$.

    For $C > 1$, $\beta > 0$, the model behaves peculiarly: all solutions consist of two rules.
    For $C=2$, $\beta = 0.5$ the ruleset consists of two adjacent, non-overlapping rules, as depicted in \autoref{fig:two_clouds_solution_bad_a}.
    The same holds for $C=2$, $\beta=1$, as shown in \autoref{fig:two_clouds_solution_bad_b}.
    These rulesets achieve a much lower objective than the rule $x_1 \geq 9.88$ found earlier.
    For $\beta=0.5,$ this can be partly attributed to the exclusion of the lower left corner in the covered region.
    More interestingly, for $\beta=1$, the choice of this ruleset can be attributed to the additive nature of the model.
    To see this, recall that points near a margin induce a loss, even if the point resides on the desired side of the covered region.
    Looking closer, constraint \eqref{eq:RM-misP} implies that the loss of a positive point decreases as the number of margins it is in increases.
    In particular, being in two or more margins can be equivalent to being far inside the covered region, which gives a loss of $0$.
    This encourages the use of more than one rule.
    From $\beta=2$ onwards, the ruleset consists of two rules, one of which is contained in the other.
    This is depicted in \autoref{fig:two_clouds_solution_bad_2}.
    The objective is again much lower than the simple rule $x_1 \geq 9.88$.
    This phenomenon can again be attributed to the additive nature of the model.

    \begin{figure}
        \centering
        \begin{subfigure}{.48\textwidth}
            \centering
            \includesvg[width=\textwidth]{gen/test_clouds/results/O_1_C_2_mw_0.5_additive.set.svg}
            \caption{}
            \label{fig:two_clouds_solution_bad_a}
        \end{subfigure}
        \hfill
        \begin{subfigure}{.48\textwidth}
            \centering
            \includesvg[width=\textwidth]{gen/test_clouds/results/O_1_C_2_mw_1_additive.set.svg}
            \caption{}
            \label{fig:two_clouds_solution_bad_b}
        \end{subfigure}
        \caption{Solutions to an instance of \autoref{ex:two_clouds} computed with the ramp loss model, configured with $C=2$, $\beta=0.5$ (a) and $C=2$, $\beta=1$ (b), respectively.}
        \label{fig:two_clouds_solution_bad}
    \end{figure}

    \begin{figure}
        \centering
        \begin{subfigure}{.48\textwidth}
            \centering
            \includesvg[width=\textwidth]{gen/test_clouds/results/O_1_C_2_mw_2_additive.set.svg}
            \caption{}
            \label{fig:two_clouds_solution_bad_2_a}
        \end{subfigure}
        \hfill
        \begin{subfigure}{.48\textwidth}
            \centering
            \includesvg[width=\textwidth]{gen/test_clouds/results/O_1_C_2_mw_3_additive.set.svg}
            \caption{}
            \label{fig:two_clouds_solution_bad_2_b}
        \end{subfigure}
        \caption{Solutions to an instance of \autoref{ex:two_clouds} computed with the ramp loss model, configured with $C=2$, $\beta=2$ (a) and $C=2$, $\beta=3$ (b), respectively.}
        \label{fig:two_clouds_solution_bad_2}
    \end{figure}

\end{example}

\begin{example}[Three clouds]\label{ex:three_clouds}
    \begin{figure}
        \includesvg[width=\textwidth]{gen/test_clouds/example_2.svg}
        \caption{Instance of \autoref{ex:three_clouds}. The example consists of a negative cloud of $100$ points, a positive cloud of $100$ points and a positive cloud of $10$ points.
        The point clouds have distributions $\text{Normal}((8, 10), 2)$, $\text{Normal}((12, 10), 2)$ and $\text{Normal}((5,15), 1)$, respectively.}
        \label{fig:three_clouds}
    \end{figure}
    Consider \autoref{ex:two_clouds} with an additional positive cloud of $10$ points, drawn from $\text{Normal}((5, 15), 1)$.
    An instance is depicted in \autoref{fig:three_clouds}.
    We again solve the ramp loss model for parameters
    \[C \in \{1,2,3\}, \quad \beta \in \{0, 0.5, 1, 2, 3, 4, 5\}, \quad D\in \{2\}.\]
    For $C = 1$, the rule $x_1 \geq 9.71$ is found for all $\beta$ up until and including $\beta = 4$.
    For $C = 1$, $\beta = 5$, the rule is $x_1 \geq 8.77$ instead.
    Notably, the margin covers every point in the positive cloud of $10$ points in this case.
    For $C = 2$ and $\beta \in \{0, 0.5, 1\},$ the ruleset
    \[
        x_1 \geq 9.71 \quad \vee \quad  \left(x_1 \leq 8.03 \; \wedge \; x_2 \geq 13.36\right)
    \]
    is found.
    This solution is depicted in \autoref{fig:three_clouds_solution_good}.
    \begin{figure}
        \includesvg[width=.8\textwidth]{gen/test_clouds/results/O_2_C_2_mw_0_additive.set.svg}
        \caption{Solution to an instance of \autoref{ex:three_clouds} computed with the ramp loss model, configured with $C=2$, $\beta=0$, $D=2$.}
        \label{fig:three_clouds_solution_good}
    \end{figure}
    For $C=2$, $\beta=2$ the resulting ruleset is different and may be considered reasonable.
    This ruleset is depicted in \autoref{fig:three_clouds_solution_bad_a}.
    For $C=2$, $\beta>2$, rulesets consist of overlapping rules and ignore the positive cloud of 10 points.
    An example is depicted in \autoref{fig:three_clouds_solution_bad_b}.
    For $C=3$, the solutions are similar to the $C=2$ case of \autoref{ex:two_clouds},
    but with an extra rule to accommodate the additional positive point cloud.
    An example is shown in \autoref{fig:three_clouds_solution_bad_2}.

    \begin{figure}
        \centering
        \begin{subfigure}{.48\textwidth}
            \centering
            \includesvg[width=\textwidth]{gen/test_clouds/results/O_2_C_2_mw_2_additive.set.svg}
            \caption{}
            \label{fig:three_clouds_solution_bad_a}
        \end{subfigure}
        \hfill
        \begin{subfigure}{.48\textwidth}
            \centering
            \includesvg[width=\textwidth]{gen/test_clouds/results/O_2_C_2_mw_3_additive.set.svg}
            \caption{}
            \label{fig:three_clouds_solution_bad_b}
        \end{subfigure}
        \caption{Solutions to an instance of \autoref{ex:three_clouds} computed with the ramp loss model, configured with $C=2$, $\beta=2$ (a) and $C=2$, $\beta=3$ (b), respectively.}
        \label{fig:three_clouds_solution_bad}
    \end{figure}

    \begin{figure}
        \includesvg[width=.8\textwidth]{gen/test_clouds/results/O_2_C_3_mw_1_additive.set.svg}
        \caption{Solution to an instance of \autoref{ex:three_clouds} computed with the ramp loss model, configured with $C=3$, $\beta=1$, $D=2$.}
        \label{fig:three_clouds_solution_bad_2}
    \end{figure}

\end{example}

In these two examples, the model behaves as desired when $\beta = 0$.
When a margin is introduced, freedom to choose rules results in undesirable rulesets.
The clearest example of this is \autoref{fig:three_clouds_solution_bad_b}, where a positive point cloud is ignored,
because overlapping rules are beneficial in terms of loss, due to additivity in the model.
From a practitioner's perspective, the problem is sometimes solved by limiting $C$ and $\beta$.
After all, for adequate choice of $C$ and $\beta$, the model produces results considered reasonable.
However, the results suggest that the model in its current form is not suitable for the problem at hand.
Therefore, we introduce a version of the model where positive points no longer benefit from being in more than one margin.
We refer to this model as the non-additive ramp loss moel.